\begin{center}
\chapter*{\size\textbf{ABSTRACT}}
\end{center}

\addcontentsline{toc}{section}{ABSTRACT}

\bigskip

\bigskip

\textbf{Title}: RE-DACT: An AI-Driven Platform for Automated Redaction of Sensitive Information in Documents

\bigskip

Digital documents often carry personally identifiable information (PII) that should not be visible when files are shared, archived, or sent for review. Manual redaction is slow and easy to get wrong, especially when the same team has to process scanned forms, identity cards, and PDFs in large numbers. This project presents RE-DACT, a web-based redaction platform for detecting and masking sensitive Indian identity information such as Aadhaar numbers, PAN numbers, payment card numbers, email addresses, and phone numbers in PDF and image files.

RE-DACT combines OCR, rule-based detection, and validation checks. Tesseract extracts text and character bounding boxes from the input document. The system then validates Aadhaar candidates with the Verhoeff algorithm and payment card candidates with the Luhn algorithm. PAN detection uses the official ten-character structure and checks the fourth character, which represents the PAN holder type. These checks reduce the false positives that would occur with plain regular-expression matching.

The processing pipeline also handles common scanning problems. It tries four orientations and two preprocessing states, giving eight OCR attempts in the worst case. If the first readable orientation produces a detection, the pipeline stops early, which keeps normal cases fast while still allowing rotated documents to be processed. In the project tests, single-page Aadhaar images with clear text were processed in 274--296 ms, while difficult cases that needed more OCR attempts took longer.

The platform provides three masking levels. Standard masking hides all but the last four characters. Partial masking hides the first four characters only, matching the current implementation. Full masking hides the entire detected value. For email addresses, phone numbers, names, and locations detected through optional NER, the system uses full masking because partial exposure is less useful and may still reveal identity.

The implementation uses FastAPI for the REST API, plain HTML/CSS/JavaScript for the browser interface, OpenCV for image operations, pdf2image for PDF conversion, and img2pdf for output assembly. It also records processing events in a JSON-lines audit log. Each entry includes a SHA-256 hash linked to the previous entry for tamper evidence during a running process. Processed files are kept only in temporary working directories, and downloadable results expire after ten minutes.

Testing on a small set of Indian identity-document samples showed reliable Aadhaar detection when the text was clear and readable. PAN, payment card, and NER detections depended more heavily on OCR quality, especially on images with complex backgrounds or stylized text. The report documents the system design, validation algorithms, implementation details, test results, limitations, and future improvements such as stronger NER, better multilingual OCR, batch processing, and containerized deployment.

\bigskip

\textbf{Keywords}: Redaction, PII, OCR, Verhoeff Algorithm, Luhn Algorithm, Tesseract, Privacy, Audit Logging, FastAPI

\bigskip

The motivation for this work stems from the practical challenge of handling sensitive Indian identity data safely in organizational workflows. Every day, government offices, educational institutions, healthcare providers, financial institutions, and private businesses receive identity proofs during onboarding, verification, claims processing, admission procedures, and compliance checks. The volume of such documents is substantial, and the manual effort required to redact them properly is correspondingly large. A small error during sharing can expose data that should have been hidden, leading to potential harm for the affected individuals and liability for the organization.

The system combines multiple detection approaches: mathematical validation for structured identifiers, pattern matching for general formats, and optional machine learning for unstructured entities. The multi-orientation OCR handling ensures that rotated and angled documents can be processed without manual pre-processing. The character-level masking provides flexibility in redaction depth appropriate for different compliance scenarios.

Testing on a small set of Indian identity-document samples showed reliable Aadhaar detection when the text was clear and readable. PAN, payment card, and NER detections depended more heavily on OCR quality, especially on images with complex backgrounds or stylized text. The report documents the system design, validation algorithms, implementation details, test results, limitations, and future improvements such as stronger NER, better multilingual OCR, batch processing, and containerized deployment.

\bigskip

The platform provides three masking levels.
